\documentclass[11pt,a4paper]{article}
\usepackage[margin=25mm]{geometry}
\usepackage{amsmath,amssymb,booktabs,microtype}
\usepackage[colorlinks,linkcolor=blue!50!black,urlcolor=blue!50!black]{hyperref}
\usepackage{xcolor}
\newcommand{\code}[1]{\texttt{\small #1}}
\newcommand{\ephi}{\bm{e}_\phi}
\usepackage{bm}
\newenvironment{keybox}{\par\medskip\noindent\rule{\linewidth}{0.4pt}\par\smallskip%
\leftskip=1em \rightskip=1em\relax}{\par\smallskip\leftskip=0pt \rightskip=0pt\noindent\rule{\linewidth}{0.4pt}\par\medskip}

\title{The prescribed floating-potential boundary condition in \code{model600}\\
\large Normalisation, sign convention, and the linear trace constraint}
\author{Derived from the JOREK source; every step checkable against a named file and line}
\date{}

\begin{document}\maketitle

\begin{abstract}
\noindent
We impose $V_p-V_{\rm wall}=\Lambda k_BT_e/e$ with $\Lambda=3$ on the plasma potential at a
material wall --- the local zero-current (floating) limit of the sheath characteristic. This note
derives, from the JOREK variables and the code's own coordinate convention, the exact linear
relation between the stream function $u$ and the electron temperature that this condition becomes,
and the two coefficients that enter the assembled matrix row. The condition is \emph{linear} in the
evolved variables, so it needs no Newton iteration, no clipping, no incidence-angle floor and no
relaxation parameter. The one genuinely delicate step is the sign, which two comments in the
current source state in opposite ways; we settle it from the coordinate basis rather than from
either comment.
\end{abstract}

\section{What is being imposed, and what is not}

For a Maxwellian plasma at a surface drawing zero net current, the sheath drops
\begin{equation}
  V_p-V_{\rm wall}=\Lambda\,\frac{k_BT_e}{e},\qquad
  \Lambda=\ln\!\sqrt{\frac{m_i}{2\pi m_e}}\approx 3 .
  \label{eq:float}
\end{equation}
This is the \emph{floating} potential: the value at which the retarded electron flux exactly
cancels the ion saturation flux. Prescribing it is a deliberate reduction of the full
Langmuir characteristic $j=j_{\rm sat}(1-e^{-X})$ to its single zero-current point.

\begin{keybox}
\textbf{Honest scope.} Eq.~\eqref{eq:float} imposes a \emph{local, zero-current} closure. It
cannot carry a net wall current, so it cannot produce thermoelectric currents between targets, and
it is not a model of a conducting vessel closing current through the wall. What it does give ---
and what the full characteristic also gives --- is a wall potential that follows $T_e$, hence a
tangential electric field and an $\bm{E}\times\bm{B}$ drift wherever $T_e$ varies along the wall.
\end{keybox}

\section{JOREK variables and the two normalisations}

The reduced-MHD ansatz uses the poloidal flux $\psi$ and the velocity stream function $u$. Two
conversions to SI are needed. With $n_0={\tt central\_density}\times10^{20}$,
$m_i={\tt central\_mass}\times m_u$ and $\rho_0=n_0m_i$,
\begin{equation}
  V_p\,[\mathrm{V}]=\frac{F_0\,u}{\sqrt{\mu_0\rho_0}},
  \qquad
  k_BT_e\,[\mathrm{J}]=\frac{T_e^{\rm JOREK}}{\mu_0 n_0}.
  \label{eq:norm}
\end{equation}
The second is JOREK's pressure normalisation ($p=\rho T$ with $\rho$ in units of $\rho_0$ and
pressure in units of $1/\mu_0$). The first is the statement $\Phi=F_0u$ in JOREK units, converted
to volts by the velocity normalisation $1/\sqrt{\mu_0\rho_0}$; its \emph{sign} is the subject of
the next section.

\section{The sign, settled from the basis and not from a comment}
\label{sec:sign}

\subsection{Why this needs care}

Two comments in the present source state opposite conventions:
\code{mod\_sheath\_bc.f90:207} says $\Phi=-F_0u$, while
\code{mod\_boundary\_conditions.f90:876} says $\Phi=+F_0u$. They cannot both be right, and the
choice flips the direction of every $\bm{E}\times\bm{B}$ drift the boundary condition drives.

The ambiguity is genuine and not resolvable from the equations alone, because $\psi$ and $u$ each
carry their own sign convention which absorbs the basis handedness. Writing the $\bm{E}\times\bm{B}$
velocity as
\begin{equation}
  \bm{v}_E=\frac{\bm{E}\times\bm{B}}{B^2},\qquad \bm{E}=-\nabla\Phi,
  \qquad
  \bm{B}\simeq\frac{F_0}{R}\ephi
  \;\Longrightarrow\;
  \bm{v}_{E,\rm pol}=-\frac{R}{F_0}\left(\nabla\Phi\times\ephi\right)
\end{equation}
and comparing with the implemented $\bm{v}_{\rm pol}=cR(\nabla u\times\ephi)$ gives $\Phi=-cF_0u$.
Everything therefore rests on the sign $c$, which depends on how $\nabla u\times\ephi$ is expanded
--- i.e.\ on the handedness of the basis.

\subsection{The basis, from three independent places in the code}

JOREK uses $(\bm{e}_R,\bm{e}_Z,\ephi)$ \emph{right-handed}, so
$\bm{e}_R\times\bm{e}_Z=+\ephi$ and $\phi$ runs clockwise viewed from above. This is not a comment;
it is visible three ways:
\begin{enumerate}\itemsep2pt
\item \code{diagnostics/jorek2\_povray.f90:102--103} and
      \code{vacuum/mod\_vacuum\_fields.f90:1172--1173} map
      $x=R\cos\phi$, $y=-R\sin\phi$. The minus fixes the handedness.
\item \code{particles/mod\_fields.f90:590--596}: \code{rot\_tmp} is the right-handed curl for
      $(R,Z,\phi)$ with scale factors $(1,1,R)$. Applied to the code's own $\bm{A}$ it returns the
      code's own $\bm{B}$ exactly; in the standard $(R,\phi,Z)$ basis it returns $-\bm{B}$ in all
      three components, so the test discriminates.
\item \code{diagnostics/new\_diag/mod\_expression.f90:1354}: $J_{\rm tor}=-z_j/R$, the third curl
      component in that basis.
\end{enumerate}

\subsection{The consequence}

In this basis $\bm{e}_R\times\ephi=-\bm{e}_Z$ and $\bm{e}_Z\times\ephi=+\bm{e}_R$, so
\begin{equation}
  \nabla u\times\ephi=\left(u_Z,\,-u_R\right).
\end{equation}
The implemented poloidal velocity is, from \code{particles/mod\_fields.f90:581--582},
\begin{equation}
  \bm{v}_{\rm pol}=\left(-R\,u_Z,\;+R\,u_R\right)=-R\left(\nabla u\times\ephi\right),
\end{equation}
i.e.\ $c=-1$, which is exactly the ansatz of H\"olzl \emph{et al.} 2021 (eq.~26). Hence
\begin{equation}
  \boxed{\;\Phi=+F_0\,u\;}
\end{equation}
and the code's $u$ equals the paper's. The $\Phi=-F_0u$ comment arose from comparing the
implemented components against the paper while implicitly using the \emph{standard}
$(R,\phi,Z)$ basis; the same components read as $-R\nabla u\times\ephi$ in JOREK's basis and
$+R\nabla u\times\ephi$ in the standard one, and the discrepancy was attributed to $u$ rather than
to the basis.

\begin{keybox}
\textbf{Consequence for this implementation.} \code{sheath\_norm} in
\code{mod\_sheath\_bc.f90} returns $a_n<0$ and must \emph{not} be reused. The floating-potential
module carries its own normalisation helper with $a_n>0$, and a unit test pins it.
\end{keybox}

\section{From the physical condition to the JOREK trace constraint}

Insert \eqref{eq:norm} into \eqref{eq:float}:
\begin{equation}
  \frac{F_0u}{\sqrt{\mu_0\rho_0}}-V_{\rm wall}
  =\frac{\Lambda}{e}\,\frac{T_e^{\rm JOREK}}{\mu_0n_0}.
\end{equation}
Solving for $u$ and introducing the two standard groups
\begin{equation}
  a_n=\frac{2eF_0\sqrt{\mu_0\rho_0}}{m_i},
  \qquad
  v_w=e\,V_{\rm wall}\,\mu_0 n_0,
\end{equation}
gives, using $\rho_0=n_0m_i$,
\begin{equation}
  \boxed{\;
  u=C_T\,T_e+C_V\,V_{\rm wall},
  \qquad
  C_T=\frac{2\Lambda}{a_n},
  \qquad
  C_V=\frac{2e\mu_0n_0}{a_n}=\frac{\sqrt{\mu_0\rho_0}}{F_0}.
  \;}
  \label{eq:constraint}
\end{equation}
Equivalently $u=2(\Lambda T_e+v_w)/a_n$. Note $a_n$ carries the sign of $F_0$, so reversing the
toroidal field reverses $u$ while leaving the reconstructed physical potential unchanged --- as it
must, since \eqref{eq:float} contains no field direction.

\paragraph{Single-temperature builds.} When \code{with\_TiTe} is false the model evolves $T=T_i+T_e$
with $T_e=T/2$, so $C_T$ is halved: $u=\Lambda T/a_n+C_VV_{\rm wall}$.

\section{Why no linearisation is required}

Eq.~\eqref{eq:constraint} is \emph{affine} in the evolved variables $(u,T_e)$ at fixed $\Lambda$.
Writing the implicit step for increments $\delta u,\delta T_e$ about the current state,
\begin{equation}
  \delta u-C_T\,\delta T_e=-\left(u^{n}-C_TT_e^{n}-C_VV_{\rm wall}\right)\equiv-g^{n},
\end{equation}
which is exact: there is no neglected higher-order term, the Jacobian entries $(1,-C_T)$ are
constants, and the residual $g^n$ vanishes identically once the constraint is satisfied. This is
the entire content of the boundary row.

\begin{keybox}
\textbf{This is the structural reason the condition is robust.} The full characteristic
$j=j_{\rm sat}(1-e^{-X})$ is nonlinear in $u$ through $\exp$, saturates at $j_{\rm sat}$, and has a
Jacobian $\partial j/\partial u\propto j_{\rm sat}\,e^{-X}$ that vanishes both in ion saturation
\emph{and} wherever $j_{\rm sat}\to0$ (grazing incidence). Every gate, slope, floor and trust region
in the current-carrying implementation exists to manage that. Eq.~\eqref{eq:constraint} has none of
those failure modes: its Jacobian is the constant $(1,-C_T)$, it never inverts a characteristic, and
it never divides by $B_n$, $g(b_n)$, $j_{\rm sat}$ or the element Jacobian.
\end{keybox}

\section{Which degrees of freedom the constraint acts on}

Because $\Lambda$ is a constant, \eqref{eq:constraint} may be differentiated along the boundary
degree of freedom by degree of freedom:
\begin{equation}
  u^{\rm val}-C_TT_e^{\rm val}=C_VV_{\rm wall},
  \qquad
  u^{\partial_\ell}-C_TT_e^{\partial_\ell}=0,
  \qquad
  u^{\partial_\ell^2}-C_TT_e^{\partial_\ell^2}=0,\;\dots
\end{equation}
where $\partial_\ell$ denotes the pure tangential derivatives carried by the trace space. The
constant $V_{\rm wall}$ appears only in the value equation, and only in the axisymmetric harmonic.

Two properties follow, and they are the reason this closure is immune to the grid pathologies that
defeated the current-carrying route:
\begin{itemize}\itemsep2pt
\item \textbf{No geometry enters.} The constraint never forms $B_n$, never divides by the element
      Jacobian $x_{\rm jac}$, and never converts logical derivatives into $(R,Z)$. The $u$ and
      $T_e$ trace DOFs are stored in the \emph{same} nodal frame, so whatever that frame is ---
      including a near-degenerate one --- its scaling cancels identically between the two terms.
\item \textbf{No normal derivative is touched.} Only the value and pure tangential DOFs are
      constrained; the normal and mixed DOFs remain governed by the PDE, exactly as for an ordinary
      Dirichlet condition.
\end{itemize}
Imposing the value alone is \emph{not} acceptable: the unconstrained Hermite tangent derivative then
generates large curvature between nodes, which $w=\Delta^\!*u$ reads directly.

\section{Numerical verification of the normalisation}

Evaluated with the constants as \code{models/constants.f90} defines them
($\mu_0=4\pi\times10^{-7}$, $e=1.602176565\times10^{-19}$,
$m_u=1.660539040\times10^{-27}$) at the AUG~38773 parameters
${\tt central\_density}=1.011088$, ${\tt central\_mass}=2.01410174369812$, $F_0=2.972306$:

\begin{center}\small
\begin{tabular}{lll}
\toprule
quantity & value & \\
\midrule
$\rho_0$                        & $3.381578331801\times10^{-7}$ & kg\,m$^{-3}$\\
$\sqrt{\mu_0\rho_0}$            & $6.518754986874\times10^{-7}$ & \\
volts per unit $u$              & $4.559622206978\times10^{6}$  & V\\
$a_n$                           & $1.856385001490\times10^{2}$  & \\
$C_T=2\Lambda/a_n$, $\Lambda=3$ & $3.232088168772\times10^{-2}$ & \\
$C_V=\sqrt{\mu_0\rho_0}/F_0$    & $2.193164158359\times10^{-7}$ & $u$ per volt\\
$T_e^{\rm JOREK}$ for $1$\,eV   & $2.035678524691\times10^{-5}$ & \\
$u$ for $T_e=1$\,eV             & $6.579492475076\times10^{-7}$ & \\
\midrule
reconstructed $V_p-V_{\rm wall}$ & $3.000000000000$ & V \\
\bottomrule
\end{tabular}
\end{center}

The round trip closes exactly, and the identity $2e\mu_0n_0/a_n=\sqrt{\mu_0\rho_0}/F_0$ holds to
all printed digits, which checks the algebra of \eqref{eq:constraint} independently.

\begin{keybox}
\textbf{Use the code's constants, not CODATA-2018.} \code{constants.f90} carries
$e=1.602176565\times10^{-19}$ and $m_u=1.660539040\times10^{-27}$, which are the 2010 values. A
unit test written with modern constants disagrees in the 8th significant figure --- close enough to
look like a rounding difference and far enough to mask a real error.
\end{keybox}

\section{What this derivation does \emph{not} establish}

\begin{itemize}\itemsep2pt
\item That the discrete problem is stable. The constraint is exact and its Jacobian is constant,
      but $u$ and $w=\Delta^\!*u$ form a mixed system and freezing boundary $w$ at a restart value
      is not a complete treatment of an evolving $u$ trace.
\item That the resulting drift is resolved. $\bm{v}_E\cdot\bm{n}\propto R\,C_T\,\partial_\ell T_e$
      is a real normal flow; whether the density advection it drives sits inside the stable
      cell-P\'eclet range of the central scheme is a separate, measurable question.
\item That a frozen $\psi$ boundary is compatible with that normal drift.
\end{itemize}
Each is an acceptance test, not an assumption, and each is measured rather than assumed in the
implementation plan.

\end{document}
